\documentclass{article}
\usepackage{amsmath,amssymb,amsthm}
\usepackage{gb4e}

\newtheorem{definition}{Definition}
\newtheorem{proposition}{Proposition}

\title{Scope Without Trees: A Toy Study of Quantifier Readings}
\author{Anonymous}

\begin{document}
\maketitle

\section{Introduction}

This note studies the readings of a handful of doubly quantified sentences.
Rather than committing to a syntax or a compositional semantics, we simply
assign each sentence one or more translations into first-order logic, as is
common in descriptive work on scope, and reason about the translations.

\begin{exe}
  \ex\label{ex:every-a} Every student read a book.
  \ex\label{ex:a-every} A student read every book.
  \ex\label{ex:some-no} Some student read no book.
  \ex\label{ex:no-every} No student read every book.
  \ex\label{ex:mary-every} Mary read every book.
  \ex\label{ex:mary-student} Mary is a student.
\end{exe}

\section{Translations}

We use a first-order language with one-place predicates $\mathit{student}$
and $\mathit{book}$, a two-place predicate $\mathit{read}$, and an individual
constant $m$ (for \textit{Mary}).

\begin{definition}[Models and entailment]\label{def:models}
A model is a standard first-order structure $M = \langle D, I\rangle$ for this
language, with $D$ non-empty. A set of formulas $\Gamma$ \emph{entails} a
formula $\varphi$ iff $\varphi$ is true in every model in which all members of
$\Gamma$ are true. Formulas are \emph{contradictory} iff no model makes them
all true.
\end{definition}

Sentence (\ref{ex:every-a}) is ambiguous between a \emph{surface} reading and
an \emph{inverse} reading:
\begin{align*}
  \text{(\ref{ex:every-a}) surface:}\quad & \forall x\,[\mathit{student}(x) \rightarrow \exists y\,[\mathit{book}(y) \wedge \mathit{read}(x,y)]]\\
  \text{(\ref{ex:every-a}) inverse:}\quad & \exists y\,[\mathit{book}(y) \wedge \forall x\,[\mathit{student}(x) \rightarrow \mathit{read}(x,y)]]
\end{align*}
Likewise for (\ref{ex:a-every}):
\begin{align*}
  \text{(\ref{ex:a-every}) surface:}\quad & \exists x\,[\mathit{student}(x) \wedge \forall y\,[\mathit{book}(y) \rightarrow \mathit{read}(x,y)]]\\
  \text{(\ref{ex:a-every}) inverse:}\quad & \forall y\,[\mathit{book}(y) \rightarrow \exists x\,[\mathit{student}(x) \wedge \mathit{read}(x,y)]]
\end{align*}
We take (\ref{ex:some-no}) and (\ref{ex:no-every}) to be unambiguous:
\begin{align*}
  \text{(\ref{ex:some-no}):}\quad & \exists x\,[\mathit{student}(x) \wedge \neg\exists y\,[\mathit{book}(y) \wedge \mathit{read}(x,y)]]\\
  \text{(\ref{ex:no-every}):}\quad & \neg\exists x\,[\mathit{student}(x) \wedge \forall y\,[\mathit{book}(y) \rightarrow \mathit{read}(x,y)]]
\end{align*}
Sentence (\ref{ex:mary-every}) is translated like the surface reading of
(\ref{ex:a-every}), with $m$ in place of the existentially bound subject, and
(\ref{ex:mary-student}) is simply $\mathit{student}(m)$.

\section{Results}

\begin{proposition}\label{prop:inv-surf-1}
The inverse reading of (\ref{ex:every-a}) entails its surface reading.
\end{proposition}

\begin{proposition}\label{prop:surf-inv-2}
The surface reading of (\ref{ex:a-every}) entails its inverse reading.
\end{proposition}

\begin{proposition}\label{prop:equiv-1}
In any model with at least one student, the two readings of
(\ref{ex:every-a}) are equivalent.
\end{proposition}

\begin{proposition}\label{prop:both-2}
Both readings of (\ref{ex:a-every}) entail that a single student read every
book.
\end{proposition}

\begin{proposition}\label{prop:no-every}
The translation of (\ref{ex:no-every}) is equivalent to
$\forall x\,[\mathit{student}(x) \rightarrow \exists y\,[\mathit{book}(y) \wedge \neg\mathit{read}(x,y)]]$:
every student failed to read at least one book.
\end{proposition}

\section{Discussion}

The proper-name sentences interact with the quantified ones as expected.
(\ref{ex:mary-every}) and (\ref{ex:mary-student}) together entail the
surface reading of (\ref{ex:a-every}), and (\ref{ex:no-every}),
(\ref{ex:mary-every}) and (\ref{ex:mary-student}) are contradictory.

Similarly, (\ref{ex:some-no}) and the surface reading of (\ref{ex:a-every})
are contradictory, since the first requires a student who read no book and
the second a student who read them all.

Finally, a model with no students verifies both readings of
(\ref{ex:every-a}) vacuously, since each reading only makes a claim about
students.

\end{document}
